% =============================================================================
% RCS Paper 5 --- Online JEPA Composite Stability (Article Format)
% Skeleton scaffold (body to be drafted across T1--T5; see PROTOCOL.md).
% =============================================================================
\documentclass[11pt,a4paper]{article}

% --- Encoding and fonts ---
\usepackage[utf8]{inputenc}
\usepackage[T1]{fontenc}
\usepackage{lmodern}

% --- Language ---
\usepackage[english]{babel}

% --- Page geometry ---
\usepackage[margin=2.5cm]{geometry}

% --- Colors (must load before hyperref for color mixing) ---
\usepackage{xcolor}

% --- Hyperlinks ---
\usepackage[colorlinks=true,linkcolor=blue!70!black,citecolor=blue!70!black,urlcolor=blue!70!black]{hyperref}

% --- Tables ---
\usepackage{booktabs}
\usepackage{array}

% --- Paragraph spacing ---
\usepackage{parskip}

% --- Micro-typography ---
\usepackage{microtype}

% --- Shared RCS preamble (math, theorems, macros) ---
% \SHAREDLATEX points at the shared LaTeX assets dir. Each paper in
% papers/pN-.../ resolves it relative to its own location; preamble.tex
% uses the macro to reference its own includes (parameters.tex, etc.).
\def\SHAREDLATEX{../../latex}
% =============================================================================
% RCS — Shared LaTeX Preamble
% Used by all papers in the Recursive Controlled Systems series
% =============================================================================

% --- Math ---
\usepackage{amsmath,amssymb,amsthm}
\usepackage{mathtools}
\usepackage{bm}              % Bold math symbols

% --- Theorem environments ---
\theoremstyle{plain}
\newtheorem{theorem}{Theorem}[section]
\newtheorem{lemma}[theorem]{Lemma}
\newtheorem{proposition}[theorem]{Proposition}
\newtheorem{corollary}[theorem]{Corollary}

\theoremstyle{definition}
\newtheorem{definition}[theorem]{Definition}
\newtheorem{example}[theorem]{Example}
\newtheorem{remark}[theorem]{Remark}
\newtheorem{assumption}[theorem]{Assumption}
\newtheorem{law}{Law}

% --- RCS notation macros ---

% Sets and spaces
\newcommand{\X}{\mathcal{X}}           % State space
\newcommand{\Y}{\mathcal{Y}}           % Observation space
\newcommand{\U}{\mathcal{U}}           % Control input space
\newcommand{\M}{\mathcal{M}}           % Mutation operator set
\newcommand{\setR}{\mathbb{R}}         % Real numbers
\newcommand{\setN}{\mathbb{N}}         % Natural numbers
\newcommand{\setZ}{\mathbb{Z}}         % Integers

% RCS 7-tuple
\newcommand{\rcs}{\Sigma}              % A Recursive Controlled System
\newcommand{\plant}{\rcs}              % Alias: plant
\newcommand{\dyn}{f}                   % Dynamics function
\newcommand{\obs}{h}                   % Observation map
\newcommand{\shield}{S}                % Safety shield
\newcommand{\ctrl}{\Pi}                % Controller

% Levels
\newcommand{\Li}[1]{L_{#1}}           % Level i: \Li{0}, \Li{1}, etc.

% Stability budget
\newcommand{\stab}{\lambda}            % Stability margin
\newcommand{\decay}{\gamma}            % Nominal decay rate
\newcommand{\adapt}{L_\theta \rho}     % Adaptation cost
\newcommand{\design}{L_d \eta}         % Design evolution cost
\newcommand{\delay}{\beta \bar{\tau}}  % Delay cost
\newcommand{\switch}{\frac{\ln \nu}{\tau_a}} % Switching cost

% Homeostatic drive
\newcommand{\drive}[1]{D_{#1}}         % Drive at level i: \drive{i}
\newcommand{\setpoint}[1]{x^*_{#1}}   % Setpoint at level i
\newcommand{\homostate}{\xi}           % Homeostatic state vector

% Lyapunov
\newcommand{\lyap}[1]{V_{#1}}         % Lyapunov function at level i
\newcommand{\dlyap}{\dot{V}}          % Time derivative of V

% EGRI
\newcommand{\egri}{\Pi_{\text{EGRI}}}  % EGRI meta-controller
\newcommand{\budget}{B}                % Budget
\newcommand{\score}{J}                 % Evaluator/score function
\newcommand{\trial}{\mathcal{T}}       % Trial record

% Agents and information state
\newcommand{\info}{I}                  % Information state (Eslami)
\newcommand{\mem}{m}                   % Memory state
\newcommand{\toolout}{z}               % Tool output
\newcommand{\interact}{r}              % Interaction signal
\newcommand{\param}{\theta}            % Adaptable parameters
\newcommand{\goal}{\zeta}              % Goal descriptor
\newcommand{\mode}{\sigma}             % Mode/strategy index
\newcommand{\arch}{c}                  % Architecture/workflow config

% Eslami delays
\newcommand{\taubar}{\bar{\tau}}       % Supremal total delay
\newcommand{\tauu}{\tau_u}             % Control computation delay
\newcommand{\tauth}{\tau_\theta}       % Adaptation update delay
\newcommand{\tauz}{\tau_z}             % Tool/retrieval delay
\newcommand{\taus}{\tau_\sigma}        % Switching decision delay
\newcommand{\tauc}{\tau_c}             % Reconfiguration delay
\newcommand{\taua}[1][]{\tau_{a\ifx\\#1\\\else,#1\fi}} % Average dwell time: \taua = tau_a, \taua[i] = tau_{a,i}

% Norms
\newcommand{\norm}[1]{\left\| #1 \right\|}
\newcommand{\abs}[1]{\left| #1 \right|}

% Operators
\DeclareMathOperator*{\argmin}{arg\,min}
\DeclareMathOperator*{\argmax}{arg\,max}


% --- Bibliography ---
\usepackage[numbers,sort&compress]{natbib}

% --- Paper-5 local macros (substrate / online JEPA) ---
% These are used only in this paper. If a macro becomes load-bearing across
% multiple papers it should be promoted to \SHAREDLATEX/preamble.tex.
\newcommand{\encoder}{s_\theta}                 % JEPA encoder
\newcommand{\predictor}{P_\varphi}              % JEPA predictor
\newcommand{\emaencoder}{s'_\theta}             % EMA target encoder
\newcommand{\head}[1]{\varphi_{H_{#1}}}         % L0 head emission map
\newcommand{\headlip}{L_H}                      % Head Lipschitz constant
\newcommand{\encoderlip}{L_\theta}              % Encoder Lipschitz constant
\newcommand{\predictorlip}{L_P}                 % Predictor Lipschitz constant
\newcommand{\onlineenc}{L_{o\theta}}            % Online encoder cost L_oθ(Δ,μ)
\newcommand{\deplag}{\Delta}                    % Deployment lag (steps)
\newcommand{\klcoef}{\mu}                       % KL regularizer coefficient
\newcommand{\headvar}{\kappa}                   % Per-level head-payload variance proxy
\newcommand{\composite}{\stab^{\mathrm{comp}}}  % Composite per-level margin

% =============================================================================
\title{Online JEPA Composite Stability\\
A Two-Time-Scale Extension of Borkar (2008)\\
for Recursive Controlled Systems with Learned Substrates}
\author{Carlos D.\ Escobar-Valbuena}
\date{May 2026 --- Scaffold Draft}

\begin{document}
\maketitle

\begin{abstract}
\emph{Scaffold placeholder.} Paper~0 of the Recursive Controlled Systems
series proves a per-level stability budget
$\stab_k = \decay_k - \adapt_k - \design_k - \delay_k - \switch_k$ under
hand-set Lyapunov functions and hand-set constants. This paper generalises
the budget to the setting in which the substrate $V_k$ itself is
\emph{learned online} by a Joint-Embedding Predictive Architecture (JEPA),
yielding two new cost terms: an \emph{online encoder cost}
$\onlineenc(\deplag, \klcoef)$ derived by extending Borkar's Ch.~6
two-time-scale stochastic approximation to discrete deployment events, and
a \emph{head cost} $\headlip\,\headvar_k$ that bounds substrate
sensitivity to head-specific payload variations. The new bound reduces to
Paper~0's Theorem~VI exactly when $\deplag\to\infty$ (frozen substrate)
and $\headlip \le 1$ (non-amplifying head); otherwise it is strictly
tighter and gates a CI-enforced production rollout via the
empirical-constants protocol of Section~\ref{sec:constants}. The final
abstract will be written in T5 once the constants table from Q3
substrate measurements is in hand. See \texttt{PROTOCOL.md} for the locked
theorem statement and proof skeleton.
\end{abstract}

\tableofcontents
\newpage

% =============================================================================
\section{Introduction}
\label{sec:introduction}
% =============================================================================
% TODO(T4): motivate moving from hand-set V_k to learned-substrate V_k;
% state the contribution (first stability bound for online JEPA substrates
% with provably-bounded deployment-lag drift); position relative to Paper~0's
% Theorem~VI and to the broader JEPA-as-substrate research program.

% =============================================================================
\section{Related Work}
\label{sec:related}
% =============================================================================
% TODO(T4): three threads:
%   1. Two-time-scale stochastic approximation: Borkar (2008), Tsitsiklis
%      \& Van Roy (1997), Konda \& Tsitsiklis (2003).
%   2. JEPA / self-supervised learning: LeCun (2022) (\citep{lecun2022path}),
%      Bardes-Ponce-LeCun (2022) VICReg, Assran et al. (2023) I-JEPA.
%   3. Stability of switched / hybrid / learned-Lyapunov systems:
%      Hespanha \& Morse (1999) (already cited in P0), Liberzon (2003),
%      Khalil (2002) -- forward to P0.

% =============================================================================
\section{RCS Recap}
\label{sec:recap}
% =============================================================================
% Short, cites Paper~0: restate the 7-tuple
% $\rcs = (\X, \Y, \U, \dyn, \obs, \shield, \ctrl)$, Paper~0's Theorem~VI
% (composite stability $\omega_N = \min_i \stab_i > 0$ under H1--H7), and
% the four hand-set constants ($\encoderlip\rho_k$, $L_d\eta_k$,
% $\beta_k\bar\tau_k$, $\ln\nu_k/\tau_a$) that Paper~5 will replace with
% \emph{learned} analogues. No new proofs here; forward all details to
% Paper~0.

% =============================================================================
\section{Theorem Statement}
\label{sec:theorem}
% =============================================================================
% TODO(T4): state the theorem verbatim from PROTOCOL.md \S1. Six assumptions
% H8a (encoder Lipschitz), H8b (predictor Lipschitz), H8c (deployment lag),
% H8d (KL-bounded drift), H8e (stability monitor), H9 (head Lipschitz emission).
% Conclusion: seven-term composite bound $\composite_k > 0$.
% Corollary: reduction to Paper~0 Theorem~VI when $\deplag\to\infty$ and
% $\headlip \le 1$.

% =============================================================================
\section{Proof}
\label{sec:proof}
% =============================================================================
% TODO(T4): five-step argument from PROTOCOL.md \S2. Lemmas land as:
%   \S\ref{sec:proof}.1 Lyapunov decomposition (T4)
%   \S\ref{sec:proof}.2 Discrete-deployment Borkar lemma (T1) ---
%                       \emph{the load-bearing technical contribution.}
%   \S\ref{sec:proof}.3 Predictor-residual bound via Lipschitz composition (T4)
%   \S\ref{sec:proof}.4 Head-Lipschitz cross-generalisation (T3)
%   \S\ref{sec:proof}.5 Sum bounds and reduction-to-P0 (T4)
% Each lemma is stated as a \texttt{lemma} environment from preamble.tex.

% =============================================================================
\section{Empirical-Constants Estimation Protocol}
\label{sec:constants}
% =============================================================================
% TODO(T5): the seven-row constants table from PROTOCOL.md \S4. CI-gated
% deployment criterion: $\hat\composite_k^{\mathrm{lower}} > 0$.
% Per-head expected $\headlip$ table (LLM, world-model, state-space, hybrid,
% behavior-cloning) from spec Section 4.4.
% The Q3 measurement script (BRO-983) populates this table with empirical
% values from at least 1000 sample pairs per Lipschitz constant.

% =============================================================================
\section{Discussion}
\label{sec:discussion}
% =============================================================================
% TODO(T4--T5): regimes where the new terms dominate; sensitivity to
% $\deplag$ and $\klcoef$ (with empirical curves from Q3); relationship to
% the JEPA-as-substrate research program (Q1--Q5); explicitly: how the
% theorem applies even when the empirical research program fails its joint
% phase gate (\S5.2 of the spec) --- the theorem's bound diagnoses the
% failure rather than being invalidated by it.

% =============================================================================
\section{Conclusion}
\label{sec:conclusion}
% =============================================================================
% TODO(T5): limitations (hand-written proof; no Lean/Coq mechanization;
% empirical figures depend on Q3 substrate landing); open questions
% (multi-head $\headlip$ amplification; non-i.i.d.\ deployment events;
% extension to L2/L3 estimators); threads to a future mechanized proof and
% to the dormant categorical-foundations Paper~5 sibling.

% =============================================================================
% Bibliography
% =============================================================================
% Scaffold: anchor the load-bearing references for Paper~5 so the
% bibliography is non-empty until T1--T5 populate the body with real
% \cite commands. Borkar 2008 is the load-bearing reference for the
% discrete-deployment two-time-scale extension (T1); Tsitsiklis \& Van
% Roy 1997 is the closest TD-learning precedent; Bishop 2006 grounds the
% KL-bound lemma (T2).
\nocite{borkar2008stochastic}
\nocite{tsitsiklis1997td}
\nocite{bishop2006prml}
\bibliographystyle{plainnat}
\bibliography{\SHAREDLATEX/references}

\end{document}
